\subsection{The complexity of approximating an instance optimum}
\label{sec:instance-hardness}

The query optimum has a finite flow characterization, proved at
the end of this appendix. Computing a close approximation within
polynomial time is harder. The following reduction starts at the
fixed row-norm bound eight, then extends to every fixed bound above
one.  The optimum \(Q^*(W)\) is defined at the start of this
appendix.

\begin{proposition}[The norm-eight reduction]
\label{prop:instance-hardness-eight}
Suppose a deterministic algorithm, polynomial in the network
description length, outputs a nonnegative rational number \(\Lambda(W)\).
Suppose there are fixed constants \(K\ge1\) and \(c\ge0\) such that,
for every width-\(n\) network with
\[
 D=\lceil\log_2n\rceil,\qquad b=20D,\qquad s\le8,
\]
it satisfies
\[
 \frac{1+Q^*(W)}{K(\log(n+2))^c}
 \le\Lambda(W)\le
 K(\log(n+2))^c(1+Q^*(W)).
\]
Then \(\mathrm P=\mathrm{NP}\).
The same approximation guarantee with probability at least \(2/3\)
from a polynomial-time randomized algorithm implies
\(\mathrm{NP}\subseteq\mathrm{BPP}\).
\end{proposition}

\begin{lemma}[Exact Boolean gates in a fixed amplitude range]
\label{lem:zero-positive-gates}
Encode false by zero and true by any number in \([1/2,1)\).
The functions
\[
 \operatorname{Or}(u_1,\ldots,u_j)
 =\tanh\left(2\sum_{i=1}^j u_i\right),\qquad 1\le j\le3,
\]
preserve this encoding and compute the Boolean disjunction.
Define
\begin{align*}
 B(u,v)={}&-\tanh(1/2+4u+4v)
       +\tanh(1/2+4u)\\
       &+\tanh(1/2+4v)-\tanh(1/2),\\
 \operatorname{And}(u,v)={}&\tanh(2B(u,v)).
\end{align*}
For encoded inputs, \(\operatorname{And}\) preserves the encoding
and computes conjunction.  The OR gate uses one layer and row norm
at most six.  The AND gate uses two layers, each with row norms at
most eight.  All biases are at most one in absolute value, and all
parameters are dyadic half-integers.
\end{lemma}

\paragraph{Proof idea.}
Encode false by exact zero and true by an interval separated from
zero. Cancellation makes the AND gate vanish exactly when an input is false,
while elementary tanh bounds keep every true output inside the encoding
interval.

\begin{proof}
For OR, the preactivation is zero when every input is zero.  If
one input is at least \(1/2\), the preactivation is at least one,
and \(\tanh1\ge1/2\).

If \(u=0\), the first and third terms of \(B\) cancel, as do
the second and fourth.  The case \(v=0\) is symmetric.  Thus a
false input makes the AND output exactly zero.  If \(u,v\ge1/2\),
monotonicity and \(\tanh x\le1\) give
\[
 B(u,v)\ge2\tanh(5/2)-1-\tanh(1/2)\ge\frac13.
\]
Here \(\tanh(1/2)\le1/2\), and
\(\tanh(5/2)\ge11/12\) follows from \(e^5\ge23\).
The exponential series gives, for example,
\(e^5\ge1+5+25/2+125/6>23\).
Also \(\tanh(2/3)\ge1/2\), since
\(e^{4/3}\ge1+4/3+8/9>3\).  Therefore
\(\operatorname{And}(u,v)\ge1/2\).  Every finite tanh output
is strictly smaller than one.

The first AND layer computes the four displayed tanh terms.
Its largest row has weights \((4,4)\), and all four biases are
\(1/2\).  The second layer has weights \((-2,2,2,-2)\)
and zero bias.  The stated parameter bounds follow.
\end{proof}

\paragraph{Proof idea.}
A formula circuit gates a balanced parity circuit. Unsatisfiable
formulas produce a constant output, while one satisfying assignment leaves
a parity signal of fixed amplitude. Any transcript missing a parity input
has zero correlation with that signal, which creates a polynomial query gap.

\begin{proof}[Proof of Proposition~\ref{prop:instance-hardness-eight}]
Start from a formula \(\varphi\) in conjunctive normal form with
at most three literals per clause, \(v\) variables, and \(k\)
clauses.  Satisfiability for these formulas
is NP-complete \cite{cook1971,karp1972}.  Let \(m=2^r\) be the least
power of two with \(m\ge\max\{v,k,2\}\); the network will have width
\(n=16m^3\).  Inputs consist of its
\(v\) assignment signs \(x\), another \(m\) signs \(y\), and
unused coordinates.  Formulas with an empty clause or with no
clauses can be decided directly.  Delete repeated literals within
each clause, and replace a clause containing complementary literals
by a constant true clause.  Pad the list of clauses
to length \(m\) with constant true clauses.

\paragraph{The formula circuit.}
At layer one, the two literal values for an assignment sign are
\[
 \tanh((1+x_i)/2),\qquad\tanh((1-x_i)/2).
\]
Exactly one is zero and the other equals \(\tanh1\ge1/2\).
At layer two, apply the OR gate to each clause's literal values.
A padded true clause is the constant \(\tanh1\), obtained from
bias one and zero weights.  Apply a balanced binary tree of AND
gates to these \(m=2^r\) clause outputs.  The final value
\(V_\varphi(x)\), at depth \(2+2r\), is zero if the assignment
does not satisfy the formula and belongs to \([1/2,1)\) otherwise.

\paragraph{The parity circuit.}
Use two outputs to encode a Boolean value: one belongs to
\([1/2,1)\) and the other is zero.  The two literal values at
layer one give such a pair for each \(y_i\).  Given encoded pairs
\((u_+,u_-)\) and \((v_+,v_-)\), their product sign is encoded by
\begin{align*}
 w_+&=\operatorname{Or}
       (\operatorname{And}(u_+,v_+),
        \operatorname{And}(u_-,v_-)),\\
 w_-&=\operatorname{Or}
       (\operatorname{And}(u_+,v_-),
        \operatorname{And}(u_-,v_+)).
\end{align*}
Only the pair matching the product sign is nonzero, and its value
is at least \(1/2\).  This gate uses three layers.  A balanced
binary tree produces \((P_+(y),P_-(y))\), encoding
\(\chi(y)=\prod_{i=1}^m y_i\), at depth \(1+3r\).

The formula circuit finishes no later than this, since
\(2+2r\le1+3r\) for \(r\ge1\).  At each remaining layer,
relay its value through \(u\mapsto\tanh(2u)\).  This relay
preserves the encoding and uses row norm two.  Denote the resulting
formula value at depth \(1+3r\) again by \(V_\varphi(x)\).

\paragraph{The output.}
Apply the AND gadget to
\((V_\varphi,P_+)\) and \((V_\varphi,P_-)\) in parallel,
obtaining \(G_+,G_-\) after two further layers.  The final neuron
is
\[
 f_\varphi(x,y)=\tanh(2(G_+(x,y)-G_-(x,y))).
\]
The total depth is
\(D=4+3r=\log_2(16m^3)=\log_2n\).
If the assignment is unsatisfying, both gated values are exactly
zero.  If it is satisfying, exactly one gated value is at least
\(1/2\), and it has the sign selected by parity.  Therefore
\begin{equation}\label{eq:SAT-parity-output}
 \varphi(x)=0\ \Longrightarrow\ f_\varphi(x,y)=0,
 \qquad
 \varphi(x)=1\ \Longrightarrow\
 \chi(y)f_\varphi(x,y)\ge\tanh1\ge\frac12.
\end{equation}

At layer one there are at most \(4m\) active neurons.  A first
layer of parallel parity-combination gates has at most \(8m\)
neurons: each pair uses four AND gates, each initially using four
neurons.  All subsequent parity layers use fewer.  The formula
circuit uses at most \(4m\) neurons in any layer, including
padding and relays.  Thus at most \(12m\) neurons are active
in a layer.  The two final AND gadgets use eight neurons in their
first layer.  Since \(12m\le16m^3=n\), zero padding gives the
required width at every layer.  All weights belong to
\(\{0,\pm1/2,\pm2,\pm4\}\), and all biases belong to
\(\{0,1/2,1\}\).  They lie on the prescribed \(2^{-20D}\)
grid.  Lemma~\ref{lem:zero-positive-gates} verifies all row norms.
The dense network description has polynomial size in \(v+k\),
and its construction takes polynomial time.

\paragraph{The cost gap.}
If \(\varphi\) is unsatisfiable, \eqref{eq:SAT-parity-output}
gives \(f_\varphi\equiv0\).  One independent fair bit samples
the output, so \(Q^*(W_\varphi)=0\).

If \(\varphi\) is satisfiable, fix one satisfying assignment
\(x^*\).  Giving this assignment to the sampler for free can
only help.  Its output sign \(Y\) satisfies, for every \(y\),
\[
 \Pr(Y=\chi(y)\mid x^*,y)
 =\frac{1+\chi(y)f_\varphi(x^*,y)}2\ge\frac34.
\]
Now draw the \(m\) parity signs independently and uniformly.
Conditional on any transcript that leaves at least one of them
unqueried, their parity is uniform.  This remains true after
including all preprocessing and online randomness in the
transcript's independent tape.  Thus, if \(Q_y\) counts the
distinct parity coordinates queried,
\[
 \Pr(Y=\chi(y))
 \le\frac12+\frac12\Pr(Q_y=m).
\]
It follows that \(\Pr(Q_y=m)\ge1/2\) and
\(\E Q_y\ge m/2\).  Averaging supplies a fixed \(y\) with
this expected cost.  This argument applies to every exact sampler
with unlimited preprocessing, so \(Q^*(W_\varphi)\ge m/2\).

Choose fixed positive integers \(K',c'\) such that the rational
quantity \(g(n)=K'(\lceil\log_2n\rceil+2)^{c'}\) bounds
\(K(\log(n+2))^c\) from above.  In the unsatisfiable case,
\(\Lambda(W_\varphi)\le g(n)\).  In the satisfiable case,
\(\Lambda(W_\varphi)\ge(1+m/2)/g(n)\).  For all sufficiently
large \(m\), \(g(16m^3)^2<1+m/2\); the two ranges are
disjoint.  Comparing the output with \(g(n)\) decides
satisfiability in polynomial time.  The finitely many smaller
formula sizes can be handled by exhaustive search.  This proves
\(\mathrm P=\mathrm{NP}\).  With a randomized approximation,
the same reduction decides satisfiability with probability at
least \(2/3\), giving \(\mathrm{NP}\subseteq\mathrm{BPP}\).
\end{proof}

\begin{theorem}[Conditional threshold for efficient instance optima]
\label{thm:instance-hardness}
Fix \(s>1\).  Suppose a deterministic polynomial-time algorithm
computes a rational value \(\Lambda(W)\) satisfying
\[
 \frac{1+Q^*(W)}{K(\log(n+2))^c}
 \le\Lambda(W)\le
 K(\log(n+2))^c(1+Q^*(W))
\]
for fixed constants \(K\ge1,c\ge0\), on every width-\(n\),
depth-\(\lceil\log_2n\rceil\) network with row norms at most
\(s\) and parameter precision \(b=20\lceil\log_2n\rceil\).
Then \(\mathrm P=\mathrm{NP}\).  A polynomial-time randomized
approximation with probability at least \(2/3\) implies
\(\mathrm{NP}\subseteq\mathrm{BPP}\).
For every fixed \(s\le1\), the constant value one approximates
\(1+Q^*(W)\) within a polylogarithmic factor on the same class.
\end{theorem}

\paragraph{Proof idea.}
For an arbitrary fixed row bound above one, insert relay layers
that restore the amplitude lost by smaller gates. Their number is a constant
depending on the row bound. Padding preserves logarithmic depth and the
polynomial gap used by the satisfiability reduction.

\begin{proof}
Fix an integer \(R\ge0\) with
\(\delta=2^{-R}\le\min\{s-1,1\}\), and put
\[
 a=1+\delta,\qquad c=\delta/16,\qquad
 K=3\cdot2^R(R+18),\qquad C=2K+3.
\]
All these quantities are fixed as the formula size increases.
We will encode false by zero and true by any value in \([c,1)\).

First consider the relay \(T(u)=\tanh(au)\).
The inequality \(\tanh v\ge v-v^3/3\) for \(v\ge0\)
follows by integrating
\(\tanh'(v)=1-\tanh^2v\ge1-v^2\).
For \(0\le u\le c\), it gives
\[
 T(u)\ge
 \left(a-\frac{a^3c^2}{3}\right)u
 \ge(1+\delta/2)u,
\]
because \(a^3c^2/3\le\delta^2/96\le\delta/2\).
Monotonicity also implies \(T(u)\ge c\) whenever \(u\ge c\).
Consequently, if
\(d=c^2/2^{14}\), then \(K\) relay layers take every input
in \([d,1)\) into \([c,1)\), and preserve zero exactly.
Indeed, until the value reaches \(c\), each relay multiplies it
by at least \(1+\delta/2\).  Moreover,
\[
 K\log(1+\delta/2)\ge K\delta/3=R+18
 \ge(R+18)\log2=\log(c/d).
\]
The inequality \(\log(1+x)\ge x/(1+x)\) gives the first bound.

Use first-layer coefficient \(\gamma=a/2\) and second-layer
coefficient \(\kappa=a/4\) in the AND construction:
\begin{align*}
 B_a(u,v)={}&-\tanh(1/2+\gamma u+\gamma v)
       +\tanh(1/2+\gamma u)\\
       &+\tanh(1/2+\gamma v)-\tanh(1/2),\\
 A_a(u,v)={}&\tanh(\kappa B_a(u,v)).
\end{align*}
The largest first-layer row norm is \(2\gamma=a\), and the
second-layer row norm is \(4\kappa=a\).
A zero input still causes exact cancellation.  If \(u,v\ge c\),
the integral identity
\[
 B_a(u,v)=\int_0^{\gamma u}\int_0^{\gamma v}
 2\tanh(1/2+y+z)\sech^2(1/2+y+z)\,dz\,dy
\]
gives an explicit positive lower bound.  Every argument belongs
to \([1/2,5/2]\).  On this interval,
\(\tanh\ge1/3\), and
\(\sech^2w\ge e^{-2w}\ge e^{-5}>1/256\).
For the last inequality, the exponential series gives \(e<3\),
so \(e^5<243<256\).  Therefore
\[
 B_a(u,v)\ge\frac{\gamma^2uv}{384}
 \ge\frac{c^2}{1536}.
\]
Also \(0\le B_a(u,v)\le2\) and \(\kappa\le1/2\).
Using \(\tanh z\ge z/2\) on \([0,1]\) gives
\[
 A_a(u,v)\ge\frac{\kappa B_a(u,v)}2
 \ge\frac{c^2}{12288}\ge d.
\]
Append \(K\) relay layers.  The resulting AND gate has depth
\(K+2\), preserves the new encoding, and has row norms at most
\(a\).

For OR with at most three inputs, use
\(\tanh(\kappa\sum_i u_i)\) followed by \(K\) relays.
Its initial row norm is at most \(3a/4\).  A true input implies
an initial output at least
\(\tanh(\kappa c)\ge\kappa c/2\ge c/8\ge d\).
Thus this OR gate has depth \(K+1\), computes disjunction in
the new encoding, and respects the norm bound.
The original literal outputs are at least \(1/2\), hence at
least \(c\), and their row norms are at most \(a\).

Repeat the formula and parity construction with these gates.
Each parity-combination level now has depth
\((K+2)+(K+1)=C\), so the parity pair is ready at depth
\(1+Cr\).  The formula value is ready at depth
\((r+1)(K+2)\), which is no larger since the difference is
\((K+1)(r-1)\ge0\).  Relay the formula value to synchronize
the two circuits.  Gating the two parity outputs by the formula
value uses \(K+2\) further layers.  Take the final output to be
\[
 f_\varphi(x,y)
 =\tanh\left(\frac a2(G_+(x,y)-G_-(x,y))\right).
\]
Its row norm is \(a\).  On a satisfying assignment its
correlation with parity is at least
\(\tanh(ac/2)\ge c/4\); on an unsatisfying assignment it
is exactly zero.

Set
\[
 n=2^{K+4}m^C,\qquad
 D=K+4+Cr=\log_2n,\qquad b=20D.
\]
The same active-width bound \(12m\) applies, since relay layers
do not increase the number of signals.  Zero padding yields the
required width \(n\).  All parameter denominators divide
\(2^{R+2}\). Since \(R+2\le20(K+4)\le b\), they lie on
the prescribed grid.
All biases remain at most one.  For fixed \(s\), \(K\) and
\(C\) are constants, so the reduction has polynomial size.

The parity transcript argument now yields
\(Q^*(W_\varphi)\ge cm/4\) in the satisfiable case and zero
in the unsatisfiable case.  This polynomial gap again exceeds
the square of any fixed polylogarithmic factor.  The approximation
consequences follow exactly as in
Proposition~\ref{prop:instance-hardness-eight}.

When \(s\le1\), Theorem~\ref{thm:upper} gives
\(Q^*(W)\le\log^{O(1)}n\).  Since \(Q^*(W)\ge0\), the
constant value one satisfies a two-sided polylogarithmic
approximation guarantee for \(1+Q^*(W)\).
\end{proof}

\begin{corollary}[Computing the multivariate score certificate]
\label{cor:score-certificate-hardness}
Let \(\mathsf I(W)\) be the first term in the supremum in
\eqref{eq:instance-certificate}, with value zero on constant
networks.  For every fixed \(s>1\), a deterministic polynomial-time
approximation to \(1+\mathsf I(W)\) within a fixed
polylogarithmic factor on the depth-\(\lceil\log_2n\rceil\),
\(b=20\lceil\log_2n\rceil\), row-norm-\(s\) class would imply
\(\mathrm P=\mathrm{NP}\).  The corresponding randomized
approximation implies \(\mathrm{NP}\subseteq\mathrm{BPP}\).
\end{corollary}

\paragraph{Proof idea.}
The balanced parity construction has the same nonzero amplitude
for every parity pattern. Fixing a satisfying assignment therefore exposes
an exact parity factor. A product distribution near the boundary makes its
multivariate score certificate polynomially large.

\begin{proof}
In the reduction, the true amplitude of a parity pair is the same
for every input pattern at a given level.  This follows by
induction: at the leaves it equals \(\tanh1\); at each balanced
combination, the correct branch applies the same AND gate to two
copies of that amplitude and the same OR gate to the result and
zero.  Relays preserve the equality.  The two final gated branches
are also identical functions of the formula value and the parity
amplitude.  Thus the network output factorizes exactly as
\[
 f_\varphi(x,y)=A_\varphi(x)\prod_{i=1}^m y_i,
\]
where \(A_\varphi(x)=0\) on an unsatisfying assignment and
\(A_\varphi(x)\ge c/4\) on a satisfying assignment, with
\(c\) as in Theorem~\ref{thm:instance-hardness}.

If the formula is unsatisfiable then \(\mathsf I(W_\varphi)=0\).
Otherwise, take product means for the assignment coordinates that
approach a satisfying assignment \(x^*\).  Give every parity
coordinate mean \(t=\sqrt{1-1/m}\), and use direction
\(v_i=1/\sqrt m\) on parity coordinates and zero on the others.
The direction is admissible in
Lemma~\ref{lem:multivariate-transcript}.  In the limit the output
mean is \(\alpha t^m\), where
\(\alpha=A_\varphi(x^*)\ge c/4\), and its directional derivative
is \(\alpha\sqrt m\,t^{m-1}\).  The denominator in the
centered information bound is at most one.  Since the product
extension and its derivatives are continuous,
\[
 \mathsf I(W_\varphi)\ge
 \alpha^2m\left(1-\frac1m\right)^{m-1}
 \ge\frac{\alpha^2m}{3}
 \ge\frac{c^2m}{48}.
\]
The middle inequality follows from
\((1+1/(m-1))^{m-1}\le e<3\).
For the fixed norm \(s\), this is a polynomial gap in the
network width.  The same comparison used for
Theorem~\ref{thm:instance-hardness} decides satisfiability.
\end{proof}

